% An ice shelf cavity in its two ventilation states.
%
% Drawn for this book. The source study's schematic is a model flow diagram;
% this one is a physical cross section, because the point here is which water
% mass reaches the ice base and not how the model is assembled.
%
% Layout rule: the cross section carries only short labels. Every sentence of
% explanation sits in a reserved band beneath it, so nothing can collide.
\begin{tikzpicture}[
  font=\footnotesize,
  note/.style={font=\scriptsize, text=black!62, align=left, inner sep=0pt},
  tag/.style={font=\scriptsize, text=black!70, inner sep=1pt, fill=white,
              fill opacity=0.75, text opacity=1},
  lbl/.style={font=\scriptsize\sffamily\bfseries, inner sep=0pt},
  pcap/.style={font=\scriptsize, text=black!75, align=left, inner sep=0pt},
  flow/.style={-{Stealth[length=4.5pt,width=3.2pt]}, line width=1.0pt},
]

\definecolor{cice}{HTML}{C9DCE6}
\definecolor{ccav}{HTML}{EAF2F6}
\definecolor{ccold}{HTML}{2E6F95}
\definecolor{cwarm}{HTML}{A8402F}
\definecolor{cbed}{HTML}{7A6A57}
\definecolor{ink}{HTML}{3A3A3A}

\def\panel#1{%
  % bedrock
  \fill[cbed!45] (0,0) -- (8.0,0) -- (8.0,0.5) -- (5.0,0.7)
       -- (2.6,0.4) -- (0,1.0) -- cycle;
  % open ocean and cavity water
  \fill[ccav] (0,1.0) -- (2.6,0.4) -- (5.0,0.7) -- (8.0,0.5)
       -- (8.0,3.6) -- (0,3.6) -- cycle;
  % grounded ice
  \fill[cice] (0,1.0) -- (0,3.6) -- (1.7,3.6) -- (1.7,1.28) -- cycle;
  % floating shelf, thinning seaward: its underside is the ice base
  \fill[cice] (1.7,1.28) -- (1.7,3.6) -- (5.6,3.6) -- (5.6,2.72) -- cycle;
  \draw[ink!55, line width=0.5pt]
    (1.7,1.28) -- (5.6,2.72) -- (5.6,3.6);
  % sea surface
  \draw[ccold!45, line width=0.6pt] (5.6,3.6) -- (8.0,3.6);
  % grounding line
  \draw[ink, line width=0.8pt] (1.7,1.28) -- (1.7,0.95);
}

% =====================================================  panel a: cold state
\begin{scope}
  \panel{}
  \node[lbl, text=ink, anchor=north west] at (0.12,3.5) {grounded ice};
  \node[lbl, text=ink] at (3.5,3.15) {ice shelf};
  \node[tag, anchor=north] at (1.7,0.92) {grounding line};
  \node[tag, anchor=south east] at (7.9,3.66) {polynya};
  \node[tag, anchor=west] at (0.15,0.22) {bedrock};
  \node[tag, text=ccold!70!black, anchor=north] at (3.4,2.18) {cavity};

  % brine rejection under the polynya
  \foreach \x in {6.1,6.5,6.9,7.3}{
    \draw[ccold, line width=0.6pt, -{Stealth[length=2.6pt,width=2.2pt]}]
      (\x,3.52) -- (\x,3.10);
  }
  % dense water descends and flows in along the bed
  \draw[flow, ccold] (6.7,2.95) -- (5.9,1.35);
  \draw[flow, ccold] (5.55,1.10) -- (2.9,0.95);
  % weak melt ticks
  \foreach \x/\y in {2.4/1.55, 3.2/1.85, 4.0/2.15, 4.8/2.45}{
    \draw[ccold!50, line width=0.4pt] (\x,\y) -- (\x,\y-0.16);
  }
  \node[lbl, text=ccold, anchor=east] at (5.45,2.62) {melt low};

  \node[pcap, anchor=north west, text width=7.9cm] at (0,-0.34)
    {(a) \textbf{Cold state.} Freezing in the polynya rejects salt, the shelf
     water becomes dense, and it sinks into the cavity. The ice base is bathed
     in water close to the freezing point, so melting is slow.};
\end{scope}

% =====================================================  panel b: warm state
\begin{scope}[xshift=9.3cm]
  \panel{}
  \node[lbl, text=ink, anchor=north west] at (0.12,3.5) {grounded ice};
  \node[lbl, text=ink] at (3.5,3.15) {ice shelf};
  \node[tag, anchor=south east] at (7.9,3.66) {weak polynya};

  % warm deep water enters along the bed and rises to the ice base
  \draw[flow, cwarm] (7.8,1.05) -- (6.0,0.95);
  \draw[flow, cwarm] (5.6,1.05) .. controls (4.4,1.35) .. (3.2,2.10);
  % strong melt ticks
  \foreach \x/\y in {2.3/1.50, 3.0/1.76, 3.7/2.02, 4.4/2.28, 5.1/2.54}{
    \draw[cwarm!80, line width=0.8pt] (\x,\y) -- (\x,\y-0.30);
  }
  \node[lbl, text=cwarm, anchor=east] at (5.45,2.62) {melt high};
  % meltwater leaves at the top of the cavity
  \draw[flow, cwarm!55] (5.0,2.86) -- (6.6,3.20);

  \node[pcap, anchor=north west, text width=7.9cm] at (0,-0.34)
    {(b) \textbf{Warm state.} Less sea ice forms, so the shelf water stays
     light and warm deep water enters instead. The meltwater it produces
     freshens the shelf further, which is what keeps the cavity warm.};
\end{scope}

\end{tikzpicture}
